
\begin{table}[H]
    \centering
\caption{Box-Typen aus der Testdatei \url{D01_V_flat_move_0001.mp4}}
\label{tab:placeholder}
    \begin{tabular}{|l|>{\raggedright\arraybackslash}p{0.85\linewidth}|}\hline
         \textbf{type}& \textbf{Beschreibung}
\\\hline
         ftyp& File Type Box: Deklariert Kompatibilität und Format (z. B. isom, 3gp4). Verifiziert die Datei als standardkonformes MP4.
\\\hline
         mdat& Media Data Box: Container für den gemultiplexten Binärstrom (Audio-/Video-Samples). Keine internen Metadaten.
\\\hline
         moov& Movie Box: Top-Level-Container für sämtliche Metadaten. Darf auf Dateiebene nur einmal vorkommen.
\\\hline
         mvhd& Movie Header Box: Enthält globale Timing-Parameter (Erstellungszeit, globale duration, timescale).
\\\hline
         trak& Track Box: Container für einen einzelnen, unabhängigen Medienstrom (z. B. Video oder Audio).
\\\hline
         tkhd& Track Header Box: Spezifiziert räumliche (Breite/Höhe in Festkommazahlen) und zeitliche Dimensionen des Tracks.
\\\hline
         mdia& Media Box: Container für die logische und zeitliche Definition der Track-Mediadaten.
\\\hline
         mdhd& Media Header Box: Definiert die trackspezifische Zeitbasis (timescale) und Dauer.
\\\hline
         hdlr& Handler Reference Box: Identifiziert den Medientyp des Tracks (z. B. vide für Video, soun für Audio).
\\ \hline
 minf&Media Information Box: Container für medienspezifische Metadaten (enthält u. a. stbl).
\\\hline
 stbl&Sample Table Box: Container für sämtliche Tabellen zur zeitlichen und physischen Lokalisierung von Samples.
\\\hline
 stsd&Sample Description Box: Definiert Codec-Spezifikationen (z. B. avc1, mp4a) und Initialisierungsparameter (SPS/PPS).
\\\hline
 stts&Time-to-Sample Box: Definiert die Dauer jedes Samples (Frames) in Ticks. Beweist VFR (Variable Framerate) oder CFR.
\\\hline
 stsz&Sample Size Box: Listet die exakte Größe (in Bytes) jedes einzelnen Samples im Track auf.
\\\hline
 stco&Chunk Offset Box: Listet die absoluten Byte-Offsets auf, bei denen Chunks (Gruppen von Samples) innerhalb der mdat beginnen.\\\hline
    \end{tabular}
    
    
\end{table}